%------------------------------------------------------------------------------
% 4.3 تصميم طبقة الرؤى الذكية (Smart Views)
% المالك: AI بالتعاون مع BE
% الحجم المستهدف: 2-3 صفحات
% الغرض: توثيق المساهمة التقنية المميزة للمشروع تصميمياً — الطبقة الوسيطة
%   بين مخطط ERP الخام ووكيل SQL، والأساس الذي يقوم عليه الأمان بالتصنيف
%   الرباعي للجداول.
%
% عناصر مطلوبة:
%   [ ] 4.3.1 بنية الطبقة: تعريف «العرض الذكي» (استعلام SQL معياري + توصيف
%       النوايا العربية المقصودة + تسميات الأعمال للمخرجات) وعلاقته بالجداول
%       المصدر، مع مثال تصميمي واحد كامل (عرض الملف الطلابي
%       v_student_profile من مستودع Bayan): الجداول المصدر، المعاملات،
%       النوايا المرتبطة، الدور المسموح
%   [ ] 4.3.2 التصنيف الأمني رباعي المستويات (tab:sv-security-tiers):
%       عام / داخلي / مقيّد بالدور / سري — لكل مستوى: من يراه، وكيف يقيّد
%       توليد SQL (رفض أم تعويض بأعم؟)، وأثره في سجل التدقيق
%   [ ] 4.3.3 دورة حياة العرض الذكي: التعريف ← المراجعة ← الإصدار ←
%       الاكتشاف التلقائي أثناء التهيئة الذكية (يرتبط بحالة الاستخدام UC-07)،
%       وكيف يتعامل النظام مع تعديل مخطط ERP المصدر (إبطال العرض المعطوب)
%
% تتبع: يخدم مباشرة حالات الاستخدام UC-01 وUC-05 وUC-07 ومتطلبات الاستعلام
%   العربي وربط مصادر البيانات بأرقامها من جدول المتطلبات الوظيفية
%------------------------------------------------------------------------------
\section{تصميم طبقة الرؤى الذكية (\EN{Smart Views})}
\label{sec:smart-views-design}
~

\subsection{بنية الطبقة وعلاقتها بمخطط \EN{ERP}}
\label{subsec:sv-structure}
~

\subsection{التصنيف الأمني رباعي المستويات للجداول}
\label{subsec:sv-security-tiers}
~

\subsection{دورة حياة العرض الذكي}
\label{subsec:sv-lifecycle}
~
